\section{Measure theory and the setting of Ergodic theory}\label{sec:measure-theory}

\subsection{Measurable spaces}\label{sec:measurable-spapces}

All along these notes, $X$ will denote a non-empty set. We shall write $2^{X}$ for the set of subsets of $X$. A \textbf{$\sigma$-algebra} on $X$ is a family $\mc{B}\subset 2^{X}$ such hat $\emptyset,X\in\mc{B}$ and given $A_{1},A_{2},\ldots\in\mc{B}$, it holds $\bigcup_{n\geq 1} A_{n}\in\mc{B}$. A \textbf{measurable space} is a pair $(X,\mc{B})$ where $X$ is a non-empty set and $\mc{B}$ is a $\sigma$-algebra on $X$.

There is a natural partial order among $\sigma$-algebras of $X$ given by set inclusion. Invoking this order, given any family $\mc{C}\subset 2^{X}$ one can define the \textbf{$\sigma$-algebra generated by $\mc{C}$}, which will be denoted by $\sigma({\mc{C}})$, as the smallest $\sigma$-algebra that contains $\mc{C}$.

When $X$ is a topological space and $\tau\subset 2^{X}$ is the topology of $X$, \ie~$\tau$ is the family of open subsets in $X$, the $\sigma$-algebra $\sigma(\tau)$ generated by $\tau$ is called the \textbf{Borel $\sigma$-algebra} of $X$ is usually denoted by $\mc{B}(X)$.


Given a measurable space $(X,\mc{B})$, and an arbitrary map $f\colon X\to Y$, one can define the \textbf{push-forward $\sigma$-algebra} $f_{\star}\mc{B}$ on $Y$ by
\begin{equation}
  \label{eq:push-forward-sigma-alg-def}
  f_{\star}\mc{B}:=\left\{E\subset Y : f^{-1}(E)\in\mc{B}\right\}.
\end{equation}

On the other hand, given another map $h\colon Z\to X$, we defined the \textbf{pull-back $\sigma$-algebra} $f^{\star}\mc{B}$ on $Z$ by
\begin{equation}
  \label{eq:pull-back-sigma-alg-def}
  f^{\star}\mc{B}:=\left\{E\subset Z : \exists A\in\B,\ f^{-1}(A)=E\right\}.
\end{equation}

Given two measurable spaces $(X,\mc{B})$ and $(Y,\mc{C})$, a map $f\colon X\to Y$ is said to be \textbf{measurable} when $\mc{C}\subset f_{\star}\mc{B}$.

\begin{exercise}
  If $(X,\mc{B})$ and $(Y,\mc{C})$ are measurable spaces and $f\colon X\to Y$, show that $f^{\star}\mc{C}\subset\mc{B}$ if and only if $\mc{C}\subset f_{\star}\mc{B}$.
\end{exercise}

A \textbf{measure} on a measurable space $(X,\mc{B})$ is map $\mu\colon\mc{B}\to[0,+\infty]$ such that $\mu(\emptyset)=0$, $\mu(X)>0$ and
\begin{displaymath}
  \mu\bigg(\bigsqcup_{n=1}^{\infty} A_{n}\bigg) = \sum_{n=1}^{\infty} \mu(A_{n}),
\end{displaymath}
for every sequence of pairwise disjoint sets $A_{1},A_{2},\ldots \in\mc{B}$.
The measure $\mu$ on $(X,\mc{B})$ is said to be \textbf{finite} when $\mu(X)<\infty$; a \textbf{probability} measure when $\mu(X)=1$; and \textbf{$\sigma$-finite} when there exists a sequence $X_{1},X_{2},\ldots\in\mc{B}$ such that $\mu(X_{n})<\infty$ for every $n\geq 1$ and $X=\bigcup_{n\geq 1} X_{n}$.

Given a measure space $(X,\mc{B},\mu)$, we say that two sets $A,B\in\mc{B}$ are \textbf{equal modulo $\mu$} when $\mu(A\triangle B)=0$, where $A\triangle B:=(A\setminus B)\cup (B\setminus A)$.

\begin{exercise}
  Let $(X,\mc{B},\mu)$ be a finite measure space. Show that the relation $\sim$ on $\mc{B}$ given by $A\sim B$ if and only if $A$ and $B$ are equal modulo $\mu$ is an equivalence relation. Then, let us consider the function $d_{\mu}\colon \mc{B}/\!\!\sim\times \mc{B}/\!\!\sim\to\R$ given by
  \begin{displaymath}
    d_{\mu}([A],[B]):=\mu(A\triangle B), \quad\forall A,B\in\B,
  \end{displaymath}
  and where $[A],[B]$ denotes the equivalence classes of $A$ and $B$, respectively. Show that the function $d_{\mu}$ is well defined and is a distance on $\mc{B}/\!\!\sim$.
\end{exercise}



A \textbf{measure space} is triple $(X,\mc{B},\mu)$, where $(X,\mc{B})$ is a measurable space and $\mu$ is a measure on it. The $\sigma$-algebra $\mc{B}_{0}^{\mu}\subset\mc{B}$ of \textbf{$\mu$-null sets} is given by
\begin{displaymath}
  \mc{B}_{0}^{\mu}:=\{N\in\mc{B} : \mu(N)=0\}.
\end{displaymath}

We say that the measure space $(X,\mc{B},\mu)$ is \textbf{complete} when for every $N\in\mc{B}_{0}^{\mu}$ and any $N'\subset N$, it holds $N\in\mc{B}$.

Every measure space $(X,\mc{B},\mu)$ admits a unique \textbf{completion:} it is a complete metric space $(X,\overline{\mc{B}},\bar\mu)$ such that
\begin{displaymath}
  \overline{\mc{B}}:=\left\{B\cup N_{1}\setminus N_{2} : B\in\mc{B},\ \exists N\in\mc{B}_{0}^{\mu},\ N_{1},N_{2}\subset N\right\},
\end{displaymath}
and
\begin{displaymath}
  \bar{\mu}(B\cup N_{1}\setminus N_{2}):=\mu(B), \quad\forall B\in\mc{B},\ \forall N\in\mc{B}_{0}^{\mu},\ \forall N_{1},N_{2}\subset N.
\end{displaymath}

So, in order to avoid unnecessary technical difficulties, from now on we will assume all our measure spaces are complete.

Given a measure space $(X,\mc{B},\mu)$ and a map $f\colon X\to Y$, we can define the \textbf{push-forward measure} $f_{\star}\mu$ on the measurable space $(Y,f_{\star}\mc{B})$ by
\begin{equation}
  \label{eq:push-forward-measure-def}
  f_{\star}\mu(A):=\mu\big(f^{-1}(A)\big), \quad\forall A\in f_{\star}\mc{B}.
\end{equation}

\begin{exercise}\label{exer:pull-back-function-push-forward-measures}
  Let $(X,\mc{B},\mu)$ be a measure space, $(Y,\mc{C})$ a measurable space and $f\colon X\to Y$ a measurable map. The it holds
  \begin{displaymath}
    \int_{X}\phi\circ f\dd\mu = \int_{Y}\phi\dd f_{\star}\mu,
  \end{displaymath}
  for any measurable function $\phi\colon Y\to [0,\infty]$.
\end{exercise}

It is important to remark that, in general, it is not possible to pulled measures back. In fact, it is only possible to do it when the corresponding map is invertible.

Given two measure spaces $(X,\mc{B},\mu)$ and $(Y,\mc{C},\nu)$, a \textbf{measure preserving map} or a \textbf{metric morphism} is a measurable map $h\colon X\to Y$ such that $h_{\star}\mu=\nu$. Observe that if we were pedantically formal, taking into account that the domain of $h_{\star}\mu$ is $h_{\star}\mc{B}\supset\mc{C}$, last equality should be written as $h_{\star}\mu\big|_{\mc{C}}=\nu$. In particular, a map $f\colon X\to X$ is said to be a \textbf{metric endomorphism} or a measure preserving map when $f$ is measurable and $f_{\star}\mu=\mu$. In such a case, we say that $\mu$ is an \textbf{$f$-invariant measure.}

A measure preserving map $h\colon (X,\mc{B},\mu)\to (Y,\mc{C},\nu)$ is said to be a \textbf{measure isomorphism} when there exist measurable sets $X'\in\mc{B}$ and $Y'\in\mc{C}$ such that $\mu(X\setminus X')=\nu(Y\setminus Y')=0$, $h\big|_{X'}\colon X'\to Y'$ is a bijection and $h^{-1}\big|_{Y'}$ is a measurable map.

Finally, we will say that a measure preserving map $f\colon (X,\mc{B},\mu)\to (X,\mc{B},\mu)$ is an \textbf{automorphism} or an \textbf{invertible measure preserving map} when $f$ is a metric isomorphism.

\subsection{Borel and Lebesgue spaces}\label{sec:Borel-Lebesgue-spaces}

Let $M$ denote a separable metric space and let $\B_{M}$ be its Borel $\sigma$-algebra, \ie~$\B_{M}$ is the smallest $\sigma$-algebra containing every open subset.

A probability space $(X,\B,\mu)$ is said to be a \textbf{Borel space} when there is a Borel probability measure $\mu_{M}$ on $(M,\B_{M})$ such that $(X,\B,\mu)$ is isomorphic to $(M,\B_{M},\mu_{M})$.

Observe that if $(X,\B,\mu)$ is Borel space, then there exists a countable family $\mc{A}\subset\B$ such that $\sigma(\mc{A})=\B$, \ie~the $\sigma$-algebra $\B$ is countably generated.

Finally, we say that the probability space $(X,\B,\mu)$ is a \textbf{Lebesgue space} when it is complete and is isomorphic to the completion of a Borel probability space. See~\cite[Appendix A. 6]{EinsiedlerWardErgThViewNumbTh} for more details about Borel and Lebesgue spaces.



\subsection{Smooth measures}\label{sec:smooth-measures}

Let $M$ and $N$ be a smooth manifolds, $\alpha$ be a $k$-form on $N$ and $f\colon M\to N$ be a $C^{1}$ map. The \textbf{pullback} of $\alpha$ by $f$ is the $k$-form $f^{\star}\alpha$ on $M$ given by
\begin{displaymath}
  f^{\star}\alpha(X_{1},X_{2},\ldots,X_{k}):= \alpha\big(Df_{x}(X_{1}),Df_{x}(X_{2}),\ldots,Df_{x}(X_{k})\big),
\end{displaymath}
for all $x\in M$, and every $X_{1},\ldots,X_{k}\in T_{x} M$.

On the other hand, given a complete smooth vector field $X$ on the manifold $M$, we can consider the flow $\Phi_{X}\colon\R\times M\to M$ induced by $X$ on $M$. We say that a Borel measure $\mu$ on $M$ is $X$-invariant (or $\Phi_{X}$-invariant) when it is $\Phi_{X}^{t}$-invariant, for every $t\in\R$, where $\Phi_{X}^{t}(z)=\Phi_{X}(z,t)$, for every $z\in M$ and any $t\in\R$.

Then, given a smooth $k$-form $\alpha$ on $M$, one defines the \emph{Lie derivative} of $\alpha$ with respect to $X$ by
\begin{equation}
  \label{eq:Lie-derivative-def}
  \mathcal{L}_{X}\alpha := \lim_{t\to 0} \frac{{\big(\Phi_{X}^{t}\big)}^{\star}\alpha-\alpha}{t}.
\end{equation}

Now, given $M$ a compact oriented smooth $d$-manifold $M$, a $d$-form $\omega$ is said to be \textbf{volume form}, when for every $x\in M$ and every positively oriented basis $\{X_{1}, X_{2},\ldots , X_{d}\}$ of the tangent space $T_{x}M$, it holds $\omega_x(X_{1},\ldots,X_{d})>0$.

Then, we can define the continuous linear functional $\Lambda_{\omega}\colon C^{\infty}(M,\R)\to\R$ by
\begin{displaymath}
  \Lambda_{\omega}(\phi):=\int \phi\omega, \quad\forall \phi\in C^{\infty}(M,\R).
\end{displaymath}
One can easily check that $\Lambda_{\omega}$ can be uniquely extended to a continuous linear functional $\overline{\Lambda_{\omega}}\colon C^{0}(M,\R)\to\R$, and by Riesz' representation theorem we conclude that there exists a finite Borel measure $\mu_{\omega}$ such that
\begin{displaymath}
  \overline{\Lambda_{\omega}}(\phi) = \int_{M}\phi\dd\mu_{\omega}, \quad\forall \phi\in C^{0}(M,\R).
\end{displaymath}

The measure $\mu_{\omega}$ is said to be a \textbf{smooth measure.}
\begin{exercise}
  Given a $C^1$-map $f\colon M\carr$, show that the smooth measure $\mu_\omega$ is $f$-invariant if and only if $f^\star\omega=\omega$.

  On the other hand, given a vector field $X$ on $M$, show that $\mu_\phi$ os $X$-invariant if and only if $\mc{L}_X\omega\equiv 0$.
\end{exercise}

\begin{exercise}
  Can we define smooth measures on non-compact oriented manifolds? Can smooth measures be finite on non-compact manifolds? Smooth measures can be defined on non-orientable manifolds?
\end{exercise}

\subsection{Some important exercises from~\cite[\S1.1.1]{VianaOliveiraFoundErgTh}}\label{sec:some-import-exerc-1.1.1}

1.1.2, 1.1.3, 1.1.4, 1.1.5.
